\documentclass[addpoints]{exam}
% \documentclass[addpoints,12pt]{exam}

\usepackage[slovene]{babel}
\usepackage{amsfonts}
\usepackage{amssymb}
\usepackage{amsmath}
\usepackage{float}
\usepackage{graphicx}
\usepackage{enumitem}
\usepackage{color}
\usepackage{eurosym}


%Komentar naslednje vrstice glede na verzijo testa -- rešitve ali prostor za odgovore:
% \printanswers

% \linespread{2}


\pointsinrightmargin
\marginbonuspointname{}


\renewcommand{\solutiontitle}{\noindent\textbf{Rešitve in točkovanje:}\par\noindent}

\colorgrids
\definecolor{GridColor}{gray}{0.7}
\setlength{\gridsize}{7mm}

\extrawidth{5mm}
\setlength{\rightpointsmargin}{1.5cm}

\begin{document}

\runningfooter{}{}{}

\begin{center}
    \fbox{\fbox{\parbox{15cm}{\centering
    Četrto preverjanje znanja \\ 1. letnik \\ 8. april 2025 \\ (Realna števila)}}}
\end{center}

\vspace{0.1in}
\makebox[\textwidth]{Ime in priimek, oddelek:\enspace\hrulefill}


% \begin{center}
%     \chqword{Naloga}
%     \chpword{Možne točke}
%     \chbpword{Dodatne točke}
%     \chsword{Dosežene točke}
%     \chtword{Skupaj}  
%     \combinedpointtable[h][questions]
%     % \combinedgradetable[h][questions]

% \end{center}

\vspace{0.1in}


\begin{questions}

    % 1. vprašanje

    \question
    Natančno izračunajte vrednost izraza.
    $$\dfrac{\sqrt{3}+1}{\sqrt{3}-1}-\dfrac{\sqrt{3}-1}{\sqrt{3}+1}+\sqrt{0.16}-\sqrt[3]{-27}+\sqrt{0.64}+\sqrt{48}-\sqrt{27}$$

    % \begin{solution}
    %     Za zapisa $b=2\cdot 3^3\cdot 5^2\cdot 7$ in $c=2^3\cdot 3^2\cdot 5^4$ po $1$ točka. 
    %     Izračun $D(a,b,c)=3^2\cdot 5^2=225$ $1$ točka, in $v(a,b,c)=2^3\cdot 3^4\cdot 5^4\cdot 7^2$ $1$ točka.
    % \end{solution}
    
    
    % 2. vprašanje
    \question
    Rešite enačbo.
        $$ \dfrac{x}{x^2-1}=\dfrac{1}{x+2}$$
    
    
    % \begin{solution}
    %     Za zapis $b=5l+1$ $1$ točka, izračun $2a+b=20k+14+5l+1$ $1$ točka, končen sklep $2a+b=5(4k+l+3)$ $1$ točka.
    % \end{solution}




    %  \newpage

    % 3. vprašanje
    \question
    Rešite sistem neenačb. Rešitev grafično upodobite in zapišite z intervalom.
    $$(x+1\leq 2x-3)\land(2x-3<0.1x+\frac{2}{5})$$
    
    % \begin{solution}[7cm]
    %     \begin{align*}
    %         1-(1+x^8)(1+x^4)(1+x^2)(1+x)(1-x)&=1-(1+x^8)(1+x^4)(1+x^2)(1-x^2)= \\
    %             &=1-(1+x^8)(1+x^4)(1-x^4)= \\
    %             &=1-(1+x^8)(1-x^8)= \\
    %             &=1-(1-x^{16})= \\
    %             &=1-1+x^{16}= \\
    %             &=x^{16}
    %     \end{align*}
    %     Za pravilno uporabo pravila razlike kvadratov $1$ točka, za pravilno upoštevan negativen predznak $1$ točka, končen rezultat $1$ točka.
    % \end{solution}


    % \begin{description}
    %     \item[a] $(x-a)^2+(y-b)^2-c^2=0$ 
    %     \item[b] $d^2(x-a)^2+f^2(y-b)^2-(df)^2=0$
    %     \item[c] $S(a,b)$
    %     \item[d] $e^2=d^2-f^2;\ d>f$
    % \end{description}
    % Krožnica: \fillin[a c][2cm]
    % \\ Elipsa:  \fillin[b c d][2cm]


        % 4. vprašanje
        \question
        Rešite sistem enačb.

        $$ \begin{aligned}
            3x+2y-4z&=2 \\ 5x-7y-5z&=-2 \\ x+5y+3z&=18
        \end{aligned} $$
    
        
        % \begin{solution}
        %     \begin{align*}
        %         1797&=1\cdot 987+610\\
        %         987&=1\cdot 610+377 \\
        %         610&=1\cdot 377+233 \\
        %         377&=1\cdot 233+144 \\
        %         233&=1\cdot 144+89 \\
        %         144&=1\cdot 89+55 \\
        %         89&=1\cdot 55+34 \\
        %         55&=1\cdot 34+21 \\
        %         34&=1\cdot 21+13 \\
        %         21&=1\cdot 13+8 \\
        %         13&=1\cdot 8+5 \\
        %         8&=1\cdot 5+3 \\
        %         5&=1\cdot 3+2 \\
        %         3&=1\cdot 2+1 \\
        %         2&=2\cdot 1+0
        %     \end{align*}
        %     Za izračun največjega skupnega delitelja $2$ točki (npr. z Evklidovim algoritmom), za pravilen sklep $1$ točka.
        % \end{solution}
    

    % \newpage
    % 5. vprašanje
    \question
        Osemnajst kuharjev pripravi obrok za veliko skupino v $40~minutah$. Po $25~minutah$ je $8$ kuharjev začelo pripravljati obrok za drugo skupino. Koliko časa bo preostanek kuharjev pripravljal obrok za prvo skupino?
    
        % \begin{solution}
        %     \begin{align*} 
        %         \frac{18}{5}\cdot 2\frac{1}{12}-\dfrac{\frac{27}{8}+1\frac{5}{6}}{\frac{19}{8}}&=\frac{18}{5}\cdot \frac{25}{12}-\dfrac{\frac{27}{8}+\frac{11}{6}}{\frac{19}{8}}= \\
        %                                                                                     &= \frac{3}{1}\cdot \frac{5}{2}-\dfrac{\frac{108}{24}+\frac{44}{24}}{\frac{19}{8}}= \\
        %                                                                                     &=\frac{15}{2}-\dfrac{\frac{152}{24}}{\frac{19}{8}}= \\
        %                                                                                     &=\frac{15}{2}-\frac{152\cdot 8}{24\cdot 19}= \\
        %                                                                                     &=\frac{15}{2}-\frac{8}{3}= \\
        %                                                                                     &=\frac{45}{6}-\frac{16}{6}= \\
        %                                                                                     &=\frac{29}{6}
        %     \end{align*}
        %     Za pravilno množenje ulomkov $1$ točka, pravilno seštevanje ulomkov $1$ točka, pravilno razreševanje dvojnih ulomkov $1$ točka, pravilno odštevanje ulomkov $1$ točka, pravilen okrajšan rezultat $1$ točka.
        % \end{solution}





    %\newpage
    % 6. vprašanje

    \question
        Natančno izračunajte. $$\left|\left|\sqrt{2}-2\right|-\left|1-\sqrt{2}\right|\right|-2$$


        % \begin{solution}
        %     \begin{align*}
        %         5(81x^2-1)(45x-5)^{-1}&= \dfrac{5(81x^2-1)}{45x-5}= \\
        %                             &= \dfrac{5(9x-1)(9x+1)}{5(9x-1)}= \\
        %                             &= 9x+1
        %     \end{align*}
        %     Pravilen zapis z ulomkom $1$ točka, razstavljanje razlike kvadratov in izpostavljanje skupnega faktorja $1$ točka, krajšanje ulomka $1$ točka, končen rezultat $1$ točka.
        % \end{solution}




    % \newpage

        % 7. vprašanje

        \question
        Rešite enačbo in opišite njen geometrijski pomen ter rešitev grafično prikažite. $$1+\left|x-7\right|=6$$


        % \begin{solution}
        %         \begin{align*}
        %             \dfrac{x+4}{x+5}-\dfrac{2x-10}{x^2-x-12}:\dfrac{x^2-25}{x-4} &= \dfrac{x+4}{x+5}-\dfrac{2(x-5)}{(x-4)(x+3)}\cdot\dfrac{x-4}{(x-5)(x+5)}= \\
        %                                                                         &= \dfrac{x+4}{x+5}-\dfrac{2}{(x+3)(x-5)}= \\
        %                                                                         &= \dfrac{(x+4)(x+3)-2}{(x+5)(x+3)}= \\
        %                                                                         &= \dfrac{x^2+7x+10}{(x+5)(x+3)}= \\
        %                                                                         &= \dfrac{(x+5)(x+2)}{(x+5)(x+3)}= \\
        %                                                                         &=\dfrac{x+2}{x+3}
        %         \end{align*}
        %         Za pravilno faktorizacijo izrazov $1$ točka, pravilno deljenje ulomkov $1$ točka, pravilno odštevanje ulomkov $1$ točka, pravilno zapisan okrajšan rezultat $1$ točka.
        %     \end{solution}
    
    


    %  \newpage

    \question
    Izračunajte vrednost izraza glede na vrednost spremenljivke $x$. $$\left|x-3\right|-\left|x+3\right|$$


    \question
    Raztopini $24~\%$ alkohola dodamo $10~kg$ $90~\%$ alkohola. Koliko kilogramov tehta mešanica, če je njena koncentracija $40~\%$?






    %  \newpage

    %  \noindent Ime in priimek: \underline{~~~~~~~~~~~~~~~~~~~~~~~~~~~~~~~~~~~~~~~~~~~~~~~~}

    %  ~\\

    %dodatno vprašanje
    % \bonusquestion[2]
    % \textit{Dodatna naloga:} Zapišite število $532$ v sedmiškem številskem sestavu.

    % \begin{solution}[7cm]
    %     \begin{align*}
    %         532 &= 7\cdot 76 + 0 \\
    %         76 &= 7\cdot 10 + 6 \\
    %         10 &= 7\cdot 1 + 3 \\
    %         1 &= 7\cdot 0 + 1
    %     \end{align*}
    %     Za pravilno zapisan algoritem $1$ točka, zapis števila v sedmiškem sistemu $532_{(10)}=1360_{(7)}$ $1$ točka.
    % \end{solution}




    % \newpage
        

    % dodatno vprašanje
    % \bonusquestion[3]
    % \textit{Dodatna naloga:} Okrajšajte ulomek: $\dfrac{b^{2x+y}+b^{x+2y}+b^{x+y+2}}{b^{2x-y}+b^x+b^{x-y+2}}$.

    % \begin{solution}[7cm]
    %     \begin{align*}
    %         \dfrac{b^{2x+y}+b^{x+2y}+b^{x+y+2}}{b^{2x-y}+b^x+b^{x-y+2}} &= \dfrac{b^{x+y}(b^x+b^y+b^2)}{b^{x-y}(b^x+b^y+b^2)}= \\
    %                                                                     &= \dfrac{b^{x+y}}{b^{x-y}}= \\
    %                                                                     &= b^{2y}
    %     \end{align*}
    %     Za pravilno izpostavljanje skupnega faktorja $1$ točka; za pravilno krajšanje $1$ točka; končen rezultat $1$ točka.
    % \end{solution}




\end{questions}



\end{document}